%% kapa: exam=1 round=34 year=2023 period=1 form=A subject=real_estate subject_ko=부동산학원론 seq=16 number=96
%% source: Q-Net 감정평가사 기출문제 (cst00309) artlSeq=5211789 file=raw/1008/5211789/2023년도 제34회 감정평가사 1차시험 1교시.pdf
%% conversion: kordoc 4.24.51 -> tools/gen1.py (auto); review=proofread-sonnet; flags=none
\begin{question}{96}
부동산투자와 위험에 관한 설명으로 옳은 것은?
\begin{choices}
  \item 상업용 부동산투자는 일반적으로 다른 상품에 비하여 초기투자비용이 많이 들며 투자비용의 회수기간이 길지만 경기침체에 민감하지 않아 투자위험이 낮다.
  \item 시장위험이란 부동산이 위치한 입지여건의 변화 때문에 발생하는 위험으로서, 부동산시장의 수요·공급과 관련된 상황의 변화와 관련되어 있다.
  \item 사업위험이란 부동산 사업자체에서 발생하는 수익성 변동의 위험을 말하며 시장위험, 입지위험, 관리·운영위험 등이 있다.
  \item 법·제도적 위험에는 소유권위험, 정부정책위험, 정치적위험, 불가항력적 위험, 유동성위험이 있다.
  \item 위험과 수익간에는 부(-)의 관계가 성립한다.
\end{choices}
\end{question}
\answer{3}
